\begin{abstract}
Agent traffic defenses must protect information as it moves between services and constrain what cover traffic executes. We systematize fifteen Agent studies, forty-four settings and twenty-two traffic-analysis anchors to connect observable signals, access rights, defense placement and execution outcomes. Two findings organize the evidence. First, tool destinations and model connections expose different representations of a task: evaluating a tool-side defense requires checking the other views available to its attacker. Second, executing cover introduces task and resource obligations that transmission overhead cannot capture. In published research-decoy evidence, the reported privacy gain and original-task utility belong to different interventions. A deployment map and transfer matrix show where established traffic-analysis guidance applies and where Agent execution adds these obligations. A controlled MCP study illustrates joint reporting of recognition, task completion, business state and cost; its deliberately adverse writes exercise the distinction between target success and collateral effects. A separate forwarding pilot leaves feedback's incremental traffic effect unresolved. We provide an evaluation procedure and a source-dependent research agenda. Published results, preprint-dependent measurements and controlled examples retain distinct evidential roles.
\end{abstract}
